\chapter{Architecture et pipeline}
\label{chap:architecture}

\section{Chaîne technique}

Le projet suit une chaîne linéaire, chaque étape consommant le résultat de la précédente sans réécrire PostgreSQL ni le fichier Excel une fois l'import effectué.

\begin{figure}[H]
\centering
\begin{tikzpicture}[
    node distance=7mm,
    box/.style={
        draw=primarycolor,
        line width=0.7pt,
        rounded corners=1pt,
        align=center,
        font=\small,
        inner sep=5pt,
        minimum width=7.4cm,
        fill=primarycolor!6
    },
    arr/.style={-{Stealth}, thick, primarycolor}
]
\node[box] (excel) {Excel\\ \texttt{marketing\_campaign\_dataset.xlsx}};
\node[box, below=of excel] (pg) {PostgreSQL\\ base \texttt{marketing\_analysis}, table unique};
\node[box, below=of pg] (sql) {SQL\\ exploration, agrégations, KPI};
\node[box, below=of sql] (py) {Python / Pandas\\ extraction lecture seule};
\node[box, below=of py] (prof) {Profiling / préparation\\ date, trimestre, CTR};
\node[box, below=of prof] (eda) {EDA\\ \texttt{python\_eda.py}};
\node[box, below=of eda] (mpl) {Matplotlib\\ 16 visualisations validées};
\node[box, below=of mpl] (st) {Streamlit\\ \texttt{dashboard.py}};
\draw[arr] (excel) -- (pg);
\draw[arr] (pg) -- (sql);
\draw[arr] (sql) -- (py);
\draw[arr] (py) -- (prof);
\draw[arr] (prof) -- (eda);
\draw[arr] (eda) -- (mpl);
\draw[arr] (mpl) -- (st);
\end{tikzpicture}
\caption{Pipeline analytique du projet (table relationnelle unique)}
\label{fig:pipeline}
\end{figure}

\section{Rôle de chaque étape}

\begin{enumerate}
\item \textbf{Excel.} Jeu source, non modifié après import.
\item \textbf{PostgreSQL.} Stockage structuré, contraintes, point de vérité pour les analyses SQL \cite{postgresql18}.
\item \textbf{SQL.} Lectures seules : volumes, moyennes, \texttt{CASE}, \texttt{HAVING}, analyses croisées et temporelles.
\item \textbf{Python / Pandas.} Extraction \texttt{SELECT *} ordonnée par \texttt{campaign\_id}, inspection des types et de la mémoire \cite{pandasdocs,sqlalchemydocs}.
\item \textbf{Profiling.} Contrôle des nulls, unicité de \texttt{campaign\_id}, conversion de la date, création des colonnes calendaires et du CTR.
\item \textbf{EDA.} Analyses A à S dans Pandas, sans graphique.
\item \textbf{Matplotlib.} Seize figures PNG, puis revue : seule l'échelle du rythme quotidien du graphique mensuel a été corrigée \cite{matplotlibdocs}.
\item \textbf{Streamlit.} Restitution filtrable des KPI et graphiques, export CSV du sous-ensemble affiché \cite{streamlitdocs}.
\end{enumerate}

\begin{infobox}[Choix d'architecture]
Le projet utilise volontairement \important{une table relationnelle unique}. Il n'y a pas de modèle en étoile, ni en flocon, ni de tables de faits et de dimensions séparées.
\end{infobox}
